\section{Ideal Simulation Results}
\label{sec:simulation-results}

The companion NumPy implementation provides reproducible checks of the paper's
core identities. It does not claim hardware realism, and exact branch
enumeration is used where sampling would obscure small differences.

\subsection{Teleportation cleanup}

Ten thousand Haar-like random single-qubit states were propagated through the
three-qubit unitary in \cref{fig:teleport-cleanup}. The maximum statevector
error after global-phase alignment was
$4.45\times 10^{-16}$; minimum fidelity was greater than
$1-1.8\times10^{-15}$. This numerically confirms the symbolic identity to
floating-point precision.

\subsection{Educational BB84 model}

One million Monte Carlo trials produced the results in
\cref{tab:simulation-results}. This simulation is a pedagogical BB84 abstraction
and is distinct from the complete Quirk-export density-matrix calculation.

\begin{table}[ht]
\centering
\caption{Ideal Monte Carlo results. ``Source access'' models duplicated
preparation controls, not the exact screenshot wire permutations.}
\label{tab:simulation-results}
\begin{tabular}{@{}lccc@{}}
\toprule
Variant & Unconditional detection & Sifted QBER & Attacker accuracy when basis matches\\
\midrule
Baseline & 0 & 0 & --\\
Intercept--resend & 0.125236 & 0.250620 & 1.000000\\
Source access & 0 & 0 & 1.000000\\
\bottomrule
\end{tabular}
\end{table}

\begin{figure}[ht]
\centering
\includegraphics[width=0.68\textwidth]{figures/bb84_detection_rates.pdf}
\caption{Ideal unconditional detection rates from one million simulated trials.
The source-access case is silent because the legitimate signal is not
intercepted.}
\label{fig:detection-rates}
\end{figure}

\subsection{Exact Quirk-export analysis}

The supplied machine-readable columns and complete output amplitudes remove the
earlier screenshot-only ambiguity. The advanced circuit matches the baseline
at level E0: their displayed $Z$-basis probability vectors agree to
$8.90\times10^{-9}$ total-variation distance. It fails at level E1: purity is
$1/2$, fidelity with the baseline is $0.4268$, and trace distance is $0.6184$.
The corresponding density matrices and derivation appear in
\cref{tab:caseii-density-metrics,fig:caseii-density-matrices}.

\subsection{Fixed-input branch-conditioned analysis}

The corrected model fixes $(v,b,e)$ and enumerates the spy's outcome $r_E$. The
principal outputs are summarized in \cref{tab:fixed-input-summary}.

\begin{table}[ht]
\centering
\caption{Information available from the retained fifth wire in the reconstructed
fixed-input model.}
\label{tab:fixed-input-summary}
\begin{tabularx}{\textwidth}{@{}lcccX@{}}
\toprule
Condition & Retained state & Value trace distance & Best value guess & Interpretation\\
\midrule
$e=b$ & $\ket{v}\!\bra{v}$ & $1$ & $1$ & Perfect value record\\
$e\neq b$ & $I/2$ & $0$ & $1/2$ & No value information locally\\
Uniform $e$, forgotten & noisy $Z$ record & $1/2$ & $3/4$ &
$0.188722$ bits of value mutual information\\
Uniform $e$, retained & conditioned pair $(e,q_4)$ & -- & -- &
$0.5$ bits of value information on average\\
\bottomrule
\end{tabularx}
\end{table}

With $v$ known and $e$ fixed, the retained state also distinguishes whether
$b=e$ with trace distance $1/2$, giving a $75\%$ optimal basis guess and
approximately $0.311278$ bits of mutual information. A single retained qubit
does not encode the four $(v,b)$ possibilities as four perfectly distinguishable
states; the result identifies a basis-conditioned value channel, not a
four-symbol classical copy.

Across every fixed branch, the maximum victim-subsystem trace distance from the
clean routed reference was $1.23\times10^{-32}$ and the minimum fidelity was
$1$. These numbers establish exactness only for the reconstructed ideal model.
The detector's original Boolean acceptance register remains unresolved, and the
public software labels the current acceptance rule provisional.

\subsection{Parametric victim-channel preservation}
\label{sec:parametric-victim}

The fixed-input result above holds for the four BB84 input states. A separate
parametric sweep extends the check to arbitrary pure-state preparations.
Alice's input qubit is set to
$|\psi(\theta,\phi)\rangle
= \cos\!\tfrac{\theta}{2}\ket{0}
+ e^{i\phi}\sin\!\tfrac{\theta}{2}\ket{1}$
for $(\theta,\phi) \in [0,\pi)\times[0,2\pi)$.

The test suite \texttt{tests/test\_parametric\_victim.py} runs two sweeps:
\begin{enumerate}
  \item A structured 60-sample grid: six values of $\theta$ by five values of
  $\phi$ for each of the two Eve bases, covering the Bloch sphere at uniform
  elevation spacing.
  \item A 200-sample pseudorandom sweep (NumPy seed 42) drawing $\theta$ and
  $\phi$ uniformly from their respective ranges.
\end{enumerate}
Across all 520 parameterized branches the maximum victim-subsystem trace
distance was below $10^{-10}$ and the minimum fidelity exceeded
$1-10^{-10}$. The Qiskit Statevector backend reproduces every branch
independently at the same tolerance, providing cross-backend confirmation.

\begin{cautionbox}[Numerical verification, not a formal channel theorem]
Sweeping over $(\theta,\phi)$ demonstrates preservation for every tested
pure-state preparation and for two independent simulation backends. It does not
constitute a formal proof that the channel is trace-preserving and
completely-positive on arbitrary inputs, including states entangled with an
external reference. A Choi-state or entangled-reference test would establish
complete channel equality and remains an open task.
The correct paper language is: \emph{``Victim preservation was numerically
verified across the tested continuous pure-state parameterization and
independently reproduced with Qiskit Statevector simulation.''}
\end{cautionbox}
